\chapter{Technology, Inevitability, and the Future}
\label{ch:technology-inevitability}

\epigraph{We shape our tools and thereafter our tools shape us.}{Marshall McLuhan (attr.)}

\noindent
The belief that technological development is autonomous, inevitable, and inherently progressive is one of the most powerful doxonomic products of modernity.
It depoliticizes technological choices by framing them as natural laws rather than social decisions.

\section{Innovation as Legitimating Discourse}
\label{sec:ch28-innovation}

\begin{definition}[Technological Determinism as Doxa]
\label{def:tech-determinism}
\index{technological determinism}
\emph{Technological determinism} is the belief that technology develops according to its own logic and that society must adapt to it.
As a doxonomic product, it serves to foreclose democratic deliberation about which technologies to develop, deploy, or constrain.
\end{definition}

The word ``innovation'' has been transformed from a neutral descriptor to a term of moral approval.
To oppose or regulate an innovation is framed as opposing progress itself.

\section{Platform Ideology}
\label{sec:ch28-platforms}

\begin{model}[Platform Doxonomics]
\label{mod:platform}
\index{platform ideology}
Social media platforms are simultaneously doxonomic infrastructure \emph{and} doxonomic actors.
They control the exposure matrix $E$ through algorithmic curation, shaping which narratives reach which audiences.
The platform's optimization function
\[
  \max_E \; \text{engagement}(E) = \max_E \sum_{j,k} e_{jk} \cdot r_{jk}
\]
where $r_{jk}$ is the revenue from exposing agent $j$ to content $k$, is not belief-neutral.
It systematically amplifies high-engagement content, which tends toward moral outrage, identity threat, and simplistic narratives \autocite{zuboff2019}.
\end{model}

\section{The Politics of Inevitability}
\label{sec:ch28-inevitability}

TINA (``There Is No Alternative'') is the purest doxonomic product: the belief that the current arrangement is not merely good but \emph{the only possibility}.
This forecloses not just opposition but imagination.

\section{Notes and References}
\label{sec:ch28-notes}

On surveillance capitalism, see \textcite{zuboff2019}.
Platform power is analyzed in \textcite{sunstein2017}.
The innovation discourse is critiqued in \textcite{mirowski2009}.
On technological determinism, see \textcite{winner1986}.
TINA and the politics of inevitability connect to \textcite{brown2015} and \textcite{streeck2016}.

\section{Exercises}

\begin{exercise}
Analyze the discourse around AI automation as a doxonomic product.
Who funds the ``AI will replace most jobs'' narrative?
What alternative framings are suppressed?
\end{exercise}

\begin{exercise}
Using the platform doxonomics model, show that engagement-maximizing algorithms reduce narrative entropy.
Under what conditions does algorithmic curation produce common-sense capture?
\end{exercise}

\begin{exercise}
Design a counter-narrative to technological determinism.
Estimate the resources needed to achieve measurable belief shift using the funding threshold from Proposition~\ref{prop:funding-threshold}.
\end{exercise}
